\documentclass[aspectratio=169,11pt]{beamer}
\usepackage[T1]{fontenc}
\usepackage{lmodern,booktabs,tabularx,array}
\definecolor{navy}{HTML}{16324F}
\definecolor{muted}{HTML}{536272}
\setbeamercolor{structure}{fg=navy}
\setbeamercolor{frametitle}{fg=navy}
\setbeamerfont{frametitle}{size=\Large,series=\bfseries}
\setbeamertemplate{navigation symbols}{}
\setbeamersize{text margin left=7mm,text margin right=7mm}
\setbeamertemplate{footline}{\leavevmode\hbox{\begin{beamercolorbox}[wd=.85\paperwidth,ht=2.6ex,dp=1.1ex,leftskip=5mm]{author in head/foot}\color{muted}\scriptsize OPALX \enspace|\enspace KV tune-shift benchmark\end{beamercolorbox}\begin{beamercolorbox}[wd=.15\paperwidth,ht=2.6ex,dp=1.1ex,rightskip=5mm]{date in head/foot}\hfill\color{muted}\scriptsize\insertframenumber/\inserttotalframenumber\end{beamercolorbox}}}
\newcolumntype{Y}{>{\raggedright\arraybackslash}X}
\newcommand{\src}[1]{\par\vfill{\color{muted}\scriptsize #1\par}}
\newcommand{\epsunit}{\ensuremath{\pi\,\mathrm{mm\,mrad}}}
\renewcommand{\arraystretch}{1.15}
\title{OPALX KV tune-shift benchmark}
\author{OPALX}
\date{26 November 2026}
\begin{document}
\begin{frame}{OPALX: benchmark parameters}
\small
\begin{tabularx}{\linewidth}{lY}
\toprule Parameter & Value used \\
\midrule
Lattice and reference radius & 16 FD cells; $R_0=12.0$ m \\
Field index and spiral angle & $k=7.630003$; $49.973388^\circ$ \\
Nominal F / D fields & $+1.650$ T / $-0.7673721$ T \\
F / D widths; F--D gap & $5.9083^\circ$ / $2.95415^\circ$; $2.38755^\circ$ \\
Fringe & Product of sigmoids; $W_0=1.17W$, steepness $16/W_0$ \\
Beam and population & Protons, $300$ MeV; $2.08\times10^{13}$ protons/bunch \\
Transverse beam & KV; equal-plane emittance cases $50,70,90$ \epsunit \\
Timestep & $0.02$ ns \\
Mesh & $128\times128$ transverse; 763 full-ring longitudinal cells \\
Longitudinal resolution & $127.9945$ cells across the initial 12 m bunch \\
\bottomrule
\end{tabularx}\par
\vspace{3mm}
\footnotesize Table I supplies the machine values, timestep and target mesh counts. Exact lattice and fringe values come from the public PyFFA configuration.
\src{Zhou et al., Phys. Rev. Accel. Beams \textbf{29}, 113301 (2026), Table I and Fig. 6; \href{https://doi.org/10.1103/q1xb-z29j}{doi:10.1103/q1xb-z29j}. PyFFA revision \texttt{2cca56b5}.}
\end{frame}

\begin{frame}{Adopted inputs and production settings}
\small
\begin{tabularx}{\linewidth}{lY}
\toprule Quantity & OPALX choice \\
\midrule
Bunch length and cut & Gaussian $\sigma_s=1.5$ m; truncate at $\pm4\sigma_s$ \\
Sampling & $131\,072$ macroparticles \\
Initial energy spread & Zero; KV beam launched with zero-current optics \\
Transverse mesh parameters & \texttt{PIPESIZEX=PIPESIZEY=0.4} m \\
Magnetic expansion & \texttt{MAX\_Y\_POWER=5}, via shared \texttt{MAXYPOWER} \\
Tracking and tune fit & Ten tracked turns; fit the last eight turns \\
Longitudinal self-fields & Disabled for this transverse comparison \\
Emittance convention & Numeric rms $\epsilon=(50,70,90)\times10^{-6}$ m rad; no extra $\pi$ \\
\bottomrule
\end{tabularx}\par
\vspace{3mm}

Cavity: 50 kV, 0 degree
dE = 0.6 MeV

\footnotesize The paper does not supply a complete KV input. These adopted beam and mesh settings limit the comparison; one mesh and timestep do not establish convergence.
\end{frame}

\begin{frame}{Maximum tune depression: OPALX and the paper}
\small
\begin{center}\setlength{\tabcolsep}{5pt}\begin{tabular}{rrrrr}
\toprule
$\epsilon$ & Paper PIC $\Delta Q_x$ & OPALX $\Delta Q_x$ & OPALX $\Delta Q_y$ & s/turn \\
\midrule
50 & -0.4913 & -0.4833 & -0.4891 & $\simeq$1162 \\
70 & -0.3505 & -0.3895 & -0.3611 & $\simeq$1162 \\
90 & -0.2636 & -0.3104 & -0.2822 & $\simeq$1162 \\
\bottomrule
\end{tabular}
\end{center}
\vspace{1mm}
\footnotesize
$\Delta Q_i=Q_i-Q_{i0}$; OPALX bare tunes $(3.2331550674,\,2.8132163504)$.
The paper uses $Q_{x0}=3.233$; its vertical coordinate is $z$.
\begin{itemize}
\item Maximum depression: minimum accepted particle tune in the central cohort $|s_{\rm initial}|<0.15$ m; no median estimator.
\item The paper does not specify its PIC core statistic or a numerical vertical PIC core tune.
\item Horizontal depression magnitude: $1.6\%$ smaller, $11.1\%$ larger and $17.8\%$ larger than the paper, respectively.
\item All cases complete ten turns without particle loss. Runtime: native \texttt{timing.dat}, two MPI ranks/two GPUs, including startup; totals round to $11\,620$ s per case.
\end{itemize}
\src{Eight-turn fits after discarding the first two turns; unchanged acceptance cuts. Extrema depend on particle count and fit quality. Paper PIC values: Zhou et al., Fig. 6. Production and analysis jobs are listed in the accompanying data.}
\end{frame}

\begin{frame}{OPALX KV footprint: $50$ \epsunit}
\centering
\includegraphics[width=\linewidth,height=.79\textheight,keepaspectratio]{figures/opalx-production-f8-50.pdf}
\src{Left: minimum accepted tune per 0.25 m initial slice; dashed curve: paper theory peak with the adopted Gaussian profile, not digitized data. Right: OPALX density; magenta: approximate paper Fig. 6 hull (CC BY 4.0), without fitting or rescaling; blue circle: bare tune.}
\end{frame}

\begin{frame}{OPALX KV footprint: $70$ \epsunit}
\centering
\includegraphics[width=\linewidth,height=.79\textheight,keepaspectratio]{figures/opalx-production-f8-70.pdf}
\src{Same slice estimator and eight-turn fit as the preceding footprint. Magenta: approximate paper Fig. 6 hull; blue circle: bare tune.}
\end{frame}

\begin{frame}{OPALX KV footprint: $90$ \epsunit}
\centering
\includegraphics[width=\linewidth,height=.79\textheight,keepaspectratio]{figures/opalx-production-f8-90.pdf}
\src{Same slice estimator and eight-turn fit as the preceding footprint. Magenta: approximate paper Fig. 6 hull; blue circle: bare tune.}
\end{frame}
\end{document}
